\chapter{svnadmin Reference---Subversion Repository
Administration}\label{svn.ref.svnadmin}

\texttt{svnadmin} is the administrative tool for monitoring and
repairing your Subversion repository. For detailed information on
repository administration, see the maintenance section for
\hyperref[svn.reposadmin.maint.tk.svnadmin]{svnadmin}.

Since \texttt{svnadmin} works via direct repository access (and thus can
only be used on the machine that holds the repository), it refers to the
repository with a path, not a URL.

Options in \texttt{svnadmin} are global, just as they are in
\texttt{svn}:

\protect\phantomsection\label{svn.ref.svnadmin.sw}{}

\begin{description}
\item[\protect\phantomsection\label{svn.ref.svnadmin.sw.bypass_hooks}{}\texttt{-\/-bypass-hooks}]
Bypass the repository hook system.
\item[\protect\phantomsection\label{svn.ref.svnadmin.sw.bypass_prop_validation}{}\texttt{-\/-bypass-prop-validation}]
When loading a dump file, disable the logic which validates property
values.
\item[\protect\phantomsection\label{svn.ref.svnadmin.sw.check_normalization}{}\texttt{-\/-check-normalization}]
When used with \texttt{svnadmin\ verify}, additionally report any names
within a single directory (or any \texttt{svn:mergeinfo} values) that are
distinct as byte strings but would collide after Unicode normalization---a
situation that can cause trouble for clients on different platforms.
(Subversion 1.9.)
\item[\protect\phantomsection\label{svn.ref.svnadmin.sw.compatible_version}{}\texttt{-\/-compatible-version}
\textless ARG\textgreater{}]
Use repository format compatible with Subversion version
\textless ARG\textgreater.
\item[\protect\phantomsection\label{svn.ref.svnadmin.sw.config_dir}{}\texttt{-\/-config-dir}
\textless DIR\textgreater{}]
Instructs Subversion to read configuration information from the
specified directory instead of the default location
(\texttt{.subversion} in the user\textquotesingle s home directory).
\item[\protect\phantomsection\label{svn.ref.svnadmin.sw.deltas}{}\texttt{-\/-deltas}]
When creating a repository dump file, specify changes in versioned
properties and file contents as deltas against their previous state.
\item[\protect\phantomsection\label{svn.ref.svnadmin.sw.exclude}{}\texttt{-\/-exclude}
\textless ARG\textgreater{}]
When used with \texttt{svnadmin\ dump}, omit from the dump any nodes whose
paths match the given prefix (or, with \texttt{-\/-pattern}, glob). May be
repeated. The effect is like an authz-based exclusion: a copy from an
excluded path is transformed into a plain add. (Subversion 1.10.)
\item[\protect\phantomsection\label{svn.ref.svnadmin.sw.file}{}\texttt{-\/-file}
(\texttt{-F}) \textless FILENAME\textgreater{}]
Uses the contents of the named file for the specified subcommand.
\item[\protect\phantomsection\label{svn.ref.svnadmin.sw.fs_type}{}\texttt{-\/-fs-type}
\textless ARG\textgreater{}]
When creating a repository, use \textless ARG\textgreater{} as the
requested filesystem type. \textless ARG\textgreater{} should be
\texttt{fsfs}.
\item[\protect\phantomsection\label{svn.ref.svnadmin.sw.force_uuid}{}\texttt{-\/-force-uuid}]
By default, when loading data into a repository that already contains
revisions, \texttt{svnadmin} will ignore the UUID from the dump stream.
This option will cause the repository\textquotesingle s UUID to be set
to the UUID from the stream.
\item[\protect\phantomsection\label{svn.ref.svnadmin.sw.ignore_uuid}{}\texttt{-\/-ignore-uuid}]
By default, when loading data into an empty repository,
\texttt{svnadmin} will set the repository\textquotesingle s UUID to the
UUID from the dump stream. This option will cause the UUID from the
stream to be ignored.
\item[\protect\phantomsection\label{svn.ref.svnadmin.sw.include}{}\texttt{-\/-include}
\textless ARG\textgreater{}]
The inverse of \texttt{-\/-exclude}: when used with
\texttt{svnadmin\ dump}, dump \emph{only} the nodes whose paths match the
given prefix (or, with \texttt{-\/-pattern}, glob). May be repeated.
(Subversion 1.10.)
\item[\protect\phantomsection\label{svn.ref.svnadmin.sw.incremental}{}\texttt{-\/-incremental}]
Dump a revision only as a diff against the previous revision, instead of
the usual fulltext.
\item[\protect\phantomsection\label{svn.ref.svnadmin.sw.keep_going}{}\texttt{-\/-keep-going}]
When used with \texttt{svnadmin\ verify}, continue checking subsequent
revisions after a corruption is found instead of stopping at the first
error, so a single run reports all detected problems. (Subversion 1.9.)
\item[\protect\phantomsection\label{svn.ref.svnadmin.sw.memory_cache_size}{}\texttt{-\/-memory-cache-size}
(\texttt{-M}) \textless ARG\textgreater{}]
Configures the size (in Megabytes) of the extra in-memory cache used to
minimize redundant operations. The default value is \texttt{16}. (This
cache is used for FSFS-backed repositories only.)
\item[\protect\phantomsection\label{svn.ref.svnadmin.sw.metadata_only}{}\texttt{-\/-metadata-only}]
When used with \texttt{svnadmin\ verify}, check only the
repository\textquotesingle s metadata and structure without reading and
verifying the full text of every representation. This is much faster but
less thorough. (Subversion 1.9.)
\item[\protect\phantomsection\label{svn.ref.svnadmin.sw.parent_dir}{}\texttt{-\/-parent-dir}
\textless DIR\textgreater{}]
When loading a dump file, root paths at \textless DIR\textgreater{}
instead of \texttt{/}.
\item[\protect\phantomsection\label{svn.ref.svnadmin.sw.pre_1.4_compatible}{}\texttt{-\/-pre-1.4-compatible}]
\emph{Deprecated}. See option \texttt{-\/-compatible-version}. When
creating a new repository, use a format that is compatible with versions
of Subversion earlier than Subversion 1.4.
\item[\protect\phantomsection\label{svn.ref.svnadmin.sw.pre_1.5_compatible}{}\texttt{-\/-pre-1.5-compatible}]
\emph{Deprecated}. See option \texttt{-\/-compatible-version}. When
creating a new repository, use a format that is compatible with versions
of Subversion earlier than Subversion 1.5.
\item[\protect\phantomsection\label{svn.ref.svnadmin.sw.pre_1.6_compatible}{}\texttt{-\/-pre-1.6-compatible}]
\emph{Deprecated}. See option \texttt{-\/-compatible-version}. When
creating a new repository, use a format that is compatible with versions
of Subversion earlier than Subversion 1.6.
\item[\protect\phantomsection\label{svn.ref.svnadmin.sw.revision}{}\texttt{-\/-revision}
(\texttt{-r}) \textless ARG\textgreater{}]
Specify a particular revision to operate on.
\item[\protect\phantomsection\label{svn.ref.svnadmin.sw.pattern}{}\texttt{-\/-pattern}]
When used with \texttt{svnadmin\ dump} together with \texttt{-\/-exclude}
or \texttt{-\/-include}, treat the supplied path prefixes as file glob
patterns (using \texttt{*}, \texttt{?}, \texttt{{[}{]}}, and
\texttt{\textbackslash{}}) rather than as literal prefixes. (Subversion
1.10.)
\item[\protect\phantomsection\label{svn.ref.svnadmin.sw.quiet}{}\texttt{-\/-quiet}
(\texttt{-q})]
Do not show normal progress---show only errors.
\item[\protect\phantomsection\label{svn.ref.svnadmin.sw.transaction}{}\texttt{-\/-transaction}
(\texttt{-t}) \textless ARG\textgreater{}]
Operate on the named uncommitted transaction instead of a committed
revision. Accepted by the \texttt{svnadmin} subcommands that can act on a
transaction, such as \texttt{verify}, \texttt{delrevprop}, and
\texttt{setrevprop}.
\item[\protect\phantomsection\label{svn.ref.svnadmin.sw.use_post_commit_hook}{}\texttt{-\/-use-post-commit-hook}]
When loading a dump file, runs the repository\textquotesingle s
post-commit hook after finalizing each newly loaded revision.
\item[\protect\phantomsection\label{svn.ref.svnadmin.sw.use_post_revprop_change_hook}{}\texttt{-\/-use-post-revprop-change-hook}]
When changing a revision property, runs the repository\textquotesingle s
post-revprop-change hook after changing the revision property.
\item[\protect\phantomsection\label{svn.ref.svnadmin.sw.use_pre_commit_hook}{}\texttt{-\/-use-pre-commit-hook}]
When loading a dump file, runs the repository\textquotesingle s
pre-commit hook before finalizing each newly loaded revision. If the
hook fails, aborts the commit and terminates the load process.
\item[\protect\phantomsection\label{svn.ref.svnadmin.sw.use_pre_revprop_change_hook}{}\texttt{-\/-use-pre-revprop-change-hook}]
When changing a revision property, runs the repository\textquotesingle s
pre-revprop-change hook before changing the revision property. If the
hook fails, aborts the modification and terminates.
\item[\protect\phantomsection\label{svn.ref.svnadmin.sw.wait}{}\texttt{-\/-wait}]
For operations which require exclusive repository access, wait until the
requisite repository lock has been obtained instead of immediately
erroring out when it cannot be.
\end{description}

\section{svnadmin
build-repcache}\label{svn.ref.svnadmin.c.build-repcache}

\emph{Build the representation cache for a repository.}

\textbf{Synopsis}

\texttt{svnadmin\ build-repcache\ REPOS\_PATH\ {[}-r\ LOWER{[}:UPPER{]}{]}}

\subsection{Description}

Add any missing entries to the repository\textquotesingle s
representation cache, the index that allows identical file and property
content to be stored only once (``representation sharing''). If a
revision range is supplied with \texttt{-\/-revision}, only those
revisions are processed; otherwise all revisions are. This is useful for
populating the cache of a repository that was created before
rep-sharing was enabled, or whose cache was lost. This command was
introduced in Subversion 1.14.

\subsection{Options}

\begin{verbatim}
--revision (-r) ARG
--quiet (-q)
\end{verbatim}

\subsection{Examples}

Populate the representation cache for an entire repository:

\begin{verbatim}
$ svnadmin build-repcache /var/svn/repos
$
\end{verbatim}

\section{svnadmin crashtest}\label{svn.ref.svnadmin.c.crashtest}

\emph{Simulate a process that crashes.}

\textbf{Synopsis}

\texttt{svnadmin\ crashtest\ REPOS\_PATH}

\subsection{Description}

Open the repository at \textless REPOS\_PATH\textgreater, then abort,
thus simulating a process that crashes while holding an open repository
handle. This is used for testing automatic repository recovery.
It\textquotesingle s unlikely that you\textquotesingle ll need to run
this command.

\subsection{Options}

None

\subsection{Examples}

\begin{verbatim}
$ svnadmin crashtest /var/svn/repos
Aborted
\end{verbatim}

Exciting, isn\textquotesingle t it?

\section{svnadmin create}\label{svn.ref.svnadmin.c.create}

\emph{Create a new, empty repository.}

\textbf{Synopsis}

\texttt{svnadmin\ create\ REPOS\_PATH}

\subsection{Description}

Create a new, empty repository at the path provided. If the provided
directory does not exist, it will be created for you.\footnote{Remember,
  \texttt{svnadmin} works only with local \emph{paths}, not \emph{URLs}.}
As of Subversion 1.2, \texttt{svnadmin} creates new repositories with
the \texttt{FSFS} filesystem backend by default.

While \texttt{svnadmin\ create} will create the base directory for a new
repository, it will not create intermediate directories. For example, if
you have an empty directory named \texttt{/var/svn}, creating
\texttt{/var/svn/repos} will work, while attempting to create
\texttt{/var/svn/subdirectory/repos} will fail with an error. Also, keep
in mind that, depending on where on your system you are creating your
repository, you might need to run \texttt{svnadmin\ create} as a user
with elevated privileges (such as the \texttt{root} user).

\subsection{Options}

\begin{verbatim}
--compatible-version ARG
--config-dir DIR
--fs-type ARG
--pre-1.4-compatible
--pre-1.5-compatible
--pre-1.6-compatible
\end{verbatim}

\subsection{Examples}

Creating a new repository is this easy:

\begin{verbatim}
$ cd /var/svn
$ svnadmin create repos
$
\end{verbatim}

\section{svnadmin delrevprop}\label{svn.ref.svnadmin.c.delrevprop}

\emph{Delete a revision property.}

\textbf{Synopsis}

\texttt{svnadmin\ delrevprop\ REPOS\_PATH\ -r\ REVISION\ NAME}

\texttt{svnadmin\ delrevprop\ REPOS\_PATH\ -t\ TXN\ NAME}

\subsection{Description}

Delete the revision property \texttt{NAME} from revision
\texttt{REVISION} (or, in the second form, from the uncommitted
transaction \texttt{TXN}). Because revision properties are not versioned,
this command \emph{irreversibly} destroys the previous value of the
property. By default the repository\textquotesingle s revision-property
change hooks are \emph{not} run; use
\texttt{-\/-use-pre-revprop-change-hook} and
\texttt{-\/-use-post-revprop-change-hook} to trigger them (for example,
to send a notification from your \texttt{post-revprop-change} hook).

\subsection{Options}

\begin{verbatim}
--revision (-r) ARG
--transaction (-t) ARG
--use-pre-revprop-change-hook
--use-post-revprop-change-hook
\end{verbatim}

\subsection{Examples}

Remove a bogus \texttt{svn:author} value accidentally set on revision 19:

\begin{verbatim}
$ svnadmin delrevprop /var/svn/repos -r 19 svn:author
$
\end{verbatim}

\section{svnadmin deltify}\label{svn.ref.svnadmin.c.deltify}

\emph{Deltify changed paths in a revision range.}

\textbf{Synopsis}

\texttt{svnadmin\ deltify\ {[}-r\ LOWER{[}:UPPER{]}{]}\ REPOS\_PATH}

\subsection{Description}

\texttt{svnadmin\ deltify} exists in current versions of Subversion only
for historical reasons. This command is deprecated and no longer needed.

It dates from a time when Subversion offered administrators greater
control over compression strategies in the repository. This turned out
to be a lot of complexity for \emph{very} little gain, and this
``feature'' was deprecated.

\subsection{Options}

\begin{verbatim}
--memory-cache-size (-M) ARG
--quiet (-q)
--revision (-r) ARG
\end{verbatim}

\section{svnadmin dump}\label{svn.ref.svnadmin.c.dump}

\emph{Dump the contents of the filesystem to \texttt{stdout}.}

\textbf{Synopsis}

\texttt{svnadmin\ dump\ REPOS\_PATH\ {[}-r\ LOWER{[}:UPPER{]}{]}\ {[}-\/-incremental{]}\ {[}-\/-deltas{]}}

\subsection{Description}

Dump the contents of the filesystem to \texttt{stdout} in a ``dump
file'' portable format, sending feedback to \texttt{stderr}. Dump
revisions \textless LOWER\textgreater{} revision through
\textless UPPER\textgreater{} revision. If no revisions are given, dump
all revision trees. If only \textless LOWER\textgreater{} is given, dump
that one revision tree. See
\hyperref[svn.reposadmin.maint.migrate]{Migrating Repository Data
Elsewhere} for a practical use.

By default, the Subversion dump stream contains a single revision (the
first revision in the requested revision range) in which every file and
directory in the repository in that revision is presented as though that
whole tree was added at once, followed by other revisions (the remainder
of the revisions in the requested range), which contain only the files
and directories that were modified in those revisions. For a modified
file, the complete full-text representation of its contents, as well as
all of its properties, are presented in the dump file; for a directory,
all of its properties are presented.

Two useful options modify the dump file generator\textquotesingle s
behavior. The first is the \texttt{-\/-incremental} option, which simply
causes that first revision in the dump stream to contain only the files
and directories modified in that revision, instead of being presented as
the addition of a new tree, and in exactly the same way that every other
revision in the dump file is presented. This is useful for generating a
relatively small dump file to be loaded into another repository that
already has the files and directories that exist in the original
repository.

The second useful option is \texttt{-\/-deltas}. This option causes
\texttt{svnadmin\ dump} to, instead of emitting full-text
representations of file contents and property lists, emit only deltas of
those items against their previous versions. This reduces (in some
cases, drastically) the size of the dump file that
\texttt{svnadmin\ dump} creates. There are, however, disadvantages to
using this option---deltified dump files are more CPU-intensive to
create and tend not to compress as well as their nondeltified
counterparts when using third-party tools such as \texttt{gzip} and
\texttt{bzip2}.

Beginning with Subversion 1.8, \texttt{svndumpfilter} can operate on
deltified dump streams. Prior to this release, \texttt{svndumpfilter}
would not work with dump streams created using \texttt{-\/-deltas}
option.

\subsection{Options}

\begin{verbatim}
--deltas
--exclude ARG
--file (-F) FILE
--include ARG
--incremental
--memory-cache-size (-M) ARG
--pattern
--quiet (-q)
--revision (-r) ARG
\end{verbatim}

\subsection{Examples}

Dump your whole repository:

\begin{verbatim}
$ svnadmin dump /var/svn/repos > full.dump
* Dumped revision 0.
* Dumped revision 1.
* Dumped revision 2.
…
\end{verbatim}

Incrementally dump a single transaction from your repository:

\begin{verbatim}
$ svnadmin dump /var/svn/repos -r 21 --incremental > incr.dump
* Dumped revision 21.
\end{verbatim}

\section{svnadmin
dump-revprops}\label{svn.ref.svnadmin.c.dump-revprops}

\emph{Dump the revision properties of a repository.}

\textbf{Synopsis}

\texttt{svnadmin\ dump-revprops\ REPOS\_PATH\ {[}-r\ LOWER{[}:UPPER{]}{]}}

\subsection{Description}

Dump only the revision \emph{properties} of the repository---not its
versioned contents---to stdout in the portable dump-file format, sending
progress feedback to stderr. If a revision range is given with
\texttt{-\/-revision}, only those revisions\textquotesingle{} properties
are dumped; otherwise the properties of all revisions are dumped. The
output can be reloaded with \texttt{svnadmin\ load-revprops}. This
command was introduced in Subversion 1.10.

\subsection{Options}

\begin{verbatim}
--revision (-r) ARG
--quiet (-q)
\end{verbatim}

\subsection{Examples}

Back up just the revision properties of a repository:

\begin{verbatim}
$ svnadmin dump-revprops /var/svn/repos > revprops.dump
* Dumped revision 0.
* Dumped revision 1.
* Dumped revision 2.
\end{verbatim}

\section{svnadmin freeze}\label{svn.ref.svnadmin.c.freeze}

\emph{Prevent commits to the repository while running an arbitary
program.}

\textbf{Synopsis}

\texttt{svnadmin\ freeze\ REPOS\_PATH\ PROGRAM\ {[}ARG...{]}}

\texttt{svnadmin\ freeze\ -\/-file\ FILENAME\ PROGRAM\ {[}ARG...{]}}

\subsection{Description}

This subcommand prevents concurrent commits to the repository
\textless REPOS\_PATH\textgreater{} (i.e. it freezes the repository)
while running \textless PROGRAM\textgreater{} with
\textless ARG\textgreater{} arguments. Clients trying to commit
concurrently will wait until the repository becomes available again. The
subcommand is intended for backup purposes so that third-party backup
tools such as \texttt{rsync} can be safely used on a live repository.

If \texttt{-\/-file} option is used, then all repositories listed in
\textless FILENAME\textgreater{} will froze. The file format is a list
of \textless REPOS\_PATH\textgreater{} separated by newlines.
Repositories freeze in the same order as they are listed in the file.

\subsection{Options}

\begin{verbatim}
--file (-F) FILENAME
\end{verbatim}

\subsection{Examples}

Freeze the repository and run \texttt{rsync} to make its backup:

\begin{verbatim}
$ svnadmin freeze /var/svn/repos -- rsync -av /var/svn/repos /backup/repos
\end{verbatim}

\section{svnadmin help (h, ?)}\label{svn.ref.svnadmin.c.help}

\emph{Help!}

\textbf{Synopsis}

\texttt{svnadmin\ help\ {[}SUBCOMMAND...{]}}

\subsection{Description}

This subcommand is useful when you\textquotesingle re trapped on a
desert island with neither a Net connection nor a copy of this book.

\section{svnadmin hotcopy}\label{svn.ref.svnadmin.c.hotcopy}

\emph{Make a hot copy of a repository.}

\textbf{Synopsis}

\texttt{svnadmin\ hotcopy\ REPOS\_PATH\ NEW\_REPOS\_PATH}

\subsection{Description}

This subcommand makes a ``hot'' backup of your repository, including all
hooks, configuration files, and, of course, database files. You can run
this command at any time and make a safe copy of the repository,
regardless of whether other processes are using the repository.

Prior to Subversion 1.8, \texttt{svnadmin\ hotcopy} always made a full
hot copy of the source repository. Beginning with Subversion 1.8 it
supports incremental copy to the existing destination copy of the same
source repository. By passing the \texttt{-\/-incremental} option to
\texttt{svnadmin\ hotcopy}, you can instruct Subversion to copy only new
revisions and revisions which have changed in size or had timestamp
modifications. The UUID of the hotcopy destination repository must match
the UUID of the hotcopy source repository. Incremental hotcopy mode is
supported for FSFS repositories only.

\subsection{Options}

\begin{verbatim}
--incremental
\end{verbatim}

\section{svnadmin info}\label{svn.ref.svnadmin.c.info}

\emph{Display information about a repository.}

\textbf{Synopsis}

\texttt{svnadmin\ info\ REPOS\_PATH}

\subsection{Description}

Print general information about the repository at \texttt{REPOS\_PATH}:
its UUID, the number of revisions, the repository and filesystem (FSFS)
format numbers, the minimum Subversion version the repository is
compatible with, and various FSFS details such as the shard size, how
many shards have been packed, and whether logical addressing is in use.
This command was introduced in Subversion 1.9.

\subsection{Options}

None

\subsection{Examples}

\begin{verbatim}
$ svnadmin info /var/svn/repos
Path: /var/svn/repos
UUID: 7d34851a-e497-4b89-b2cf-639b91bda4ba
Revisions: 2
Repository Format: 5
Compatible With Version: 1.10.0
Repository Capability: mergeinfo
Filesystem Type: fsfs
Filesystem Format: 8
FSFS Sharded: yes
FSFS Shard Size: 1000
FSFS Shards Packed: 0/0
FSFS Logical Addressing: yes
Configuration File: /var/svn/repos/db/fsfs.conf
$
\end{verbatim}

\section{svnadmin load}\label{svn.ref.svnadmin.c.load}

\emph{Read a repository dump stream from \texttt{stdin}.}

\textbf{Synopsis}

\texttt{svnadmin\ load\ REPOS\_PATH\ {[}-r\ LOWER{[}:UPPER{]}{]}}

\subsection{Description}

Read a repository dump stream from \texttt{stdin}, committing new
revisions into the repository\textquotesingle s filesystem. Progress
feedback is sent to \texttt{stdout}. If no revisions are given, read and
commit all revisions. But if \texttt{-\/-revision} (\texttt{-r}) is
provided, read and commit revisions \textless LOWER\textgreater{}
revision through \textless UPPER\textgreater{} revision only. If only
\textless LOWER\textgreater{} is given, load that one revision.

Prior to Subversion 1.8, \texttt{svnadmin\ load} was limited to reading
all revisions that the dump stream contains, but now
\texttt{svnadmin\ load} accepts \texttt{-\/-revision} (\texttt{-r})
option that acts as a filter for dump stream revisions. This allows you
to incrementally load only a range of revisions from a single dump
stream making some repository maintenance and reorganization tasks much
easier.

\subsection{Options}

\begin{verbatim}
--bypass-prop-validation
--force-uuid
--ignore-uuid
--memory-cache-size (-M) ARG
--parent-dir DIR
--quiet (-q)
--revision (-r) ARG
--use-post-commit-hook
--use-pre-commit-hook
\end{verbatim}

\subsection{Examples}

This shows the beginning of loading a repository from a backup file
(made, of course, with \texttt{svnadmin\ dump}):

\begin{verbatim}
$ svnadmin load /var/svn/restored < repos-backup
<<< Started new txn, based on original revision 1
     * adding path : test ... done.
     * adding path : test/a ... done.
…
\end{verbatim}

Or if you want to load into a subdirectory:

\begin{verbatim}
$ svnadmin load --parent-dir new/subdir/for/project \
                /var/svn/restored < repos-backup
<<< Started new txn, based on original revision 1
     * adding path : test ... done.
     * adding path : test/a ... done.
…
\end{verbatim}

Newer versions of Subversion have grown more strict regarding the format
of the values of Subversion\textquotesingle s own built-in properties.
Of course, properties created with older versions of Subversion
wouldn\textquotesingle t have benefitted from that strictness, and as
such might be improperly formatted. Dump streams carry property values
as-is, so using Subversion 1.8 to load dump streams created from
repositories with ill-formatted property values will, by default,
trigger a validation error. There are several workaround for this
problem. First, you can manually repair the problematic property values
in the source repository and recreate the dump stream. Or, you can
manually tweak the dump stream itself to fix those property values.
Finally, if you\textquotesingle d rather not deal with the problem right
now, use the \texttt{-\/-bypass-prop-validation} option with
\texttt{svnadmin\ load}.

\section{svnadmin
load-revprops}\label{svn.ref.svnadmin.c.load-revprops}

\emph{Load revision properties into a repository.}

\textbf{Synopsis}

\texttt{svnadmin\ load-revprops\ REPOS\_PATH}

\subsection{Description}

Read a dump-file-formatted stream from stdin and set the revision
properties in the repository\textquotesingle s filesystem---the reverse
of \texttt{svnadmin\ dump-revprops}. Revisions present in the stream but
not in the repository cause an error. If \texttt{-\/-revision} is given,
only the matching revisions from the stream are loaded. This command was
introduced in Subversion 1.10.

\subsection{Options}

\begin{verbatim}
--revision (-r) ARG
--quiet (-q)
--bypass-prop-validation
--force-uuid
--normalize-props
--no-flush-to-disk
\end{verbatim}

\subsection{Examples}

Restore the revision properties dumped earlier with
\texttt{svnadmin\ dump-revprops}:

\begin{verbatim}
$ svnadmin load-revprops /var/svn/repos < revprops.dump
Properties set on revision 0.
Properties set on revision 1.
Properties set on revision 2.
\end{verbatim}

\section{svnadmin lock}\label{svn.ref.svnadmin.c.lock}

\emph{Lock path in the repository directly.}

\textbf{Synopsis}

\texttt{svnadmin\ lock\ REPOS\_PATH\ PATH-IN-REPOS\ USERNAME\ FILE\ {[}TOKEN{]}}

\subsection{Description}

Lock \textless PATH-IN-REPOS\textgreater{} in the repository, assigning
ownership of the lock to \textless USERNAME\textgreater{} and using the
contents of \textless FILE\textgreater{} as comments associated with the
created lock. If provided, use \textless TOKEN\textgreater{} as lock
token.

\subsection{Options}

\begin{verbatim}
--bypass-hooks
\end{verbatim}

\section{svnadmin lslocks}\label{svn.ref.svnadmin.c.lslocks}

\emph{Print descriptions of all locks.}

\textbf{Synopsis}

\texttt{svnadmin\ lslocks\ REPOS\_PATH\ {[}PATH-IN-REPOS{]}}

\subsection{Description}

Print descriptions of all locks in repository
\textless REPOS\_PATH\textgreater{} underneath the path
\textless PATH-IN-REPOS\textgreater. If
\textless PATH-IN-REPOS\textgreater{} is not provided, it defaults to
the root directory of the repository.

\subsection{Options}

None

\subsection{Examples}

This lists the one locked file in the repository at
\texttt{/var/svn/repos}:

\begin{verbatim}
$ svnadmin lslocks /var/svn/repos
Path: /tree.jpg
UUID Token: opaquelocktoken:ab00ddf0-6afb-0310-9cd0-dda813329753
Owner: harry
Created: 2005-07-08 17:27:36 -0500 (Fri, 08 Jul 2005)
Expires: 
Comment (1 line):
Rework the uppermost branches on the bald cypress in the foreground.
\end{verbatim}

\section{svnadmin lstxns}\label{svn.ref.svnadmin.c.lstxns}

\emph{Print the names of all uncommitted transactions.}

\textbf{Synopsis}

\texttt{svnadmin\ lstxns\ REPOS\_PATH}

\subsection{Description}

Print the names of all uncommitted transactions. See
\hyperref[svn.reposadmin.maint.diskspace.deadtxns]{Removing dead
transactions} for information on how uncommitted transactions are
created and what you should do with them.

\subsection{Examples}

List all outstanding transactions in a repository:

\begin{verbatim}
$ svnadmin lstxns /var/svn/repos/ 
1w
1x
\end{verbatim}

\section{svnadmin pack}\label{svn.ref.svnadmin.c.pack}

\emph{Possibly compact the repository into a more efficient storage
model.}

\textbf{Synopsis}

\texttt{svnadmin\ pack\ REPOS\_PATH}

\subsection{Description}

See \hyperref[svn.reposadmin.maint.diskspace.fsfspacking]{Packing FSFS
filesystems} for more information.

\subsection{Options}

None

\section{svnadmin recover}\label{svn.ref.svnadmin.c.recover}

\emph{Bring a repository back into a consistent state. In addition, if
\texttt{repos/conf/passwd} does not exist, it will create a default
password file.}

\textbf{Synopsis}

\texttt{svnadmin\ recover\ REPOS\_PATH}

\subsection{Description}

Run this command if you get an error indicating that your repository
needs to be recovered.

\subsection{Options}

\begin{verbatim}
--wait
\end{verbatim}

\subsection{Examples}

Recover a hung repository:

\begin{verbatim}
$ svnadmin recover /var/svn/repos/ 
Repository lock acquired.
Please wait; recovering the repository may take some time...

Recovery completed.
The latest repos revision is 34.
\end{verbatim}

Recovering the database requires an exclusive lock on the repository.
(This is a ``database lock''; see the sidebar
\hyperref[svn.advanced.locking.meanings]{The Many Meanings of
``Lock''}.) If another process is accessing the repository, then
\texttt{svnadmin\ recover} will error:

\begin{verbatim}
$ svnadmin recover /var/svn/repos
svn: E165000: Failed to get exclusive repository access; perhaps another proce
ss such as httpd, svnserve or svn has it open?
$
\end{verbatim}

The \texttt{-\/-wait} option, however, will cause
\texttt{svnadmin\ recover} to wait indefinitely for other processes to
disconnect:

\begin{verbatim}
$ svnadmin recover /var/svn/repos --wait
Waiting on repository lock; perhaps another process has it open?

### time goes by…

Repository lock acquired.
Please wait; recovering the repository may take some time...

Recovery completed.
The latest repos revision is 34.
\end{verbatim}

\section{svnadmin rev-size}\label{svn.ref.svnadmin.c.rev-size}

\emph{Measure the size of a revision.}

\textbf{Synopsis}

\texttt{svnadmin\ rev-size\ REPOS\_PATH\ -r\ REVISION}

\subsection{Description}

Print the total size, in bytes, of the on-disk representation of revision
\texttt{REVISION}. The reported size includes the revision\textquotesingle s
properties but excludes FSFS indexes. With \texttt{-\/-quiet}, only the
number and a newline are printed, which is convenient for scripting. This
command was introduced in Subversion 1.13.

\subsection{Options}

\begin{verbatim}
--revision (-r) ARG
--quiet (-q)
\end{verbatim}

\subsection{Examples}

\begin{verbatim}
$ svnadmin rev-size /var/svn/repos -r 2
         800 bytes in revision 2
$ svnadmin rev-size /var/svn/repos -r 2 --quiet
800
\end{verbatim}

\section{svnadmin rmlocks}\label{svn.ref.svnadmin.c.rmlocks}

\emph{Unconditionally remove one or more locks from a repository.}

\textbf{Synopsis}

\texttt{svnadmin\ rmlocks\ REPOS\_PATH\ LOCKED\_PATH...}

\subsection{Description}

Remove one or more locks from each \textless LOCKED\_PATH\textgreater.

\subsection{Options}

None

\subsection{Examples}

This deletes the locks on \texttt{tree.jpg} and \texttt{house.jpg} in
the repository at \texttt{/var/svn/repos}:

\begin{verbatim}
$ svnadmin rmlocks /var/svn/repos tree.jpg house.jpg
Removed lock on '/tree.jpg.
Removed lock on '/house.jpg.
\end{verbatim}

\section{svnadmin rmtxns}\label{svn.ref.svnadmin.c.rmtxns}

\emph{Delete transactions from a repository.}

\textbf{Synopsis}

\texttt{svnadmin\ rmtxns\ REPOS\_PATH\ TXN\_NAME...}

\subsection{Description}

Delete outstanding transactions from a repository. This is covered in
detail in \hyperref[svn.reposadmin.maint.diskspace.deadtxns]{Removing
dead transactions}.

\subsection{Options}

\begin{verbatim}
--quiet (-q)
\end{verbatim}

\subsection{Examples}

Remove named transactions:

\begin{verbatim}
$ svnadmin rmtxns /var/svn/repos/ 1w 1x
\end{verbatim}

Fortunately, the output of \texttt{lstxns} works great as the input for
\texttt{rmtxns}:

\begin{verbatim}
$ svnadmin rmtxns /var/svn/repos/  `svnadmin lstxns /var/svn/repos/`
\end{verbatim}

This removes all uncommitted transactions from your repository.

\section{svnadmin setlog}\label{svn.ref.svnadmin.c.setlog}

\emph{Set the log message on a revision.}

\textbf{Synopsis}

\texttt{svnadmin\ setlog\ REPOS\_PATH\ -r\ REVISION\ FILE}

\subsection{Description}

Set the log message on revision \textless REVISION\textgreater{} to the
contents of \textless FILE\textgreater.

This is similar to using \texttt{svn\ propset} with the
\texttt{-\/-revprop} option to set the \texttt{svn:log} property on a
revision, except that you can also use the option
\texttt{-\/-bypass-hooks} to avoid running any pre- or post-commit
hooks, which is useful if the modification of revision properties has
not been enabled in the pre-revprop-change hook.

Revision properties are not under version control, so this command will
permanently overwrite the previous log message.

\subsection{Options}

\begin{verbatim}
--bypass-hooks
--revision (-r) ARG
\end{verbatim}

\subsection{Examples}

Set the log message for revision 19 to the contents of the file
\texttt{msg}:

\begin{verbatim}
$ svnadmin setlog /var/svn/repos/ -r 19 msg
\end{verbatim}

\section{svnadmin setrevprop}\label{svn.ref.svnadmin.c.setrevprop}

\emph{Set a property on a revision.}

\textbf{Synopsis}

\texttt{svnadmin\ setrevprop\ REPOS\_PATH\ -r\ REVISION\ NAME\ FILE}

\subsection{Description}

Set the property \textless NAME\textgreater{} on revision
\textless REVISION\textgreater{} to the contents of
\textless FILE\textgreater. Use \texttt{-\/-use-pre-revprop-change-hook}
or \texttt{-\/-use-post-revprop-change-hook} to trigger the revision
property-related hooks (e.g., if you want an email notification sent
from your post-revprop-change hook).

\subsection{Options}

\begin{verbatim}
--revision (-r) ARG
--use-post-revprop-change-hook
--use-pre-revprop-change-hook
\end{verbatim}

\subsection{Examples}

The following sets the revision property \texttt{repository-photo} to
the contents of the file \texttt{sandwich.png}:

\begin{verbatim}
$ svnadmin setrevprop /var/svn/repos -r 0 repository-photo sandwich.png
\end{verbatim}

As you can see, \texttt{svnadmin\ setrevprop} has no output upon
success.

\section{svnadmin setuuid}\label{svn.ref.svnadmin.c.setuuid}

\emph{Reset the repository UUID.}

\textbf{Synopsis}

\texttt{svnadmin\ setuuid\ REPOS\_PATH\ {[}NEW\_UUID{]}}

\subsection{Description}

Reset the repository UUID for the repository located at
\textless REPOS\_PATH\textgreater. If \textless NEW\_UUID\textgreater{}
is provided, use that as the new repository UUID; otherwise, generate a
brand-new UUID for the repository.

\subsection{Options}

None

\subsection{Examples}

If you\textquotesingle ve \texttt{svnsync}ed \texttt{/var/svn/repos} to
\texttt{/var/svn/repos-new} and intend to use \texttt{repos-new} as your
canonical repository, you may want to change the UUID for
\texttt{repos-new} to the UUID of \texttt{repos} so that your users
don\textquotesingle t have to check out a new working copy to
accommodate the change:

\begin{verbatim}
$ svnadmin setuuid /var/svn/repos-new 2109a8dd-854f-0410-ad31-d604008985ab
\end{verbatim}

As you can see, \texttt{svnadmin\ setuuid} has no output upon success.

\section{svnadmin unlock}\label{svn.ref.svnadmin.c.unlock}

\emph{Unlock path in the repository directly.}

\textbf{Synopsis}

\texttt{svnadmin\ unlock\ REPOS\_PATH\ LOCKED\_PATH\ USERNAME\ TOKEN}

\subsection{Description}

Unlock \textless LOCKED\_PATH\textgreater{} in the repository (as
\textless USERNAME\textgreater) after verifying that the token
associated with the lock matches \textless TOKEN\textgreater.

\subsection{Options}

\begin{verbatim}
--bypass-hooks
\end{verbatim}

\section{svnadmin upgrade}\label{svn.ref.svnadmin.c.upgrade}

\emph{Upgrade a repository to the latest supported schema version.}

\textbf{Synopsis}

\texttt{svnadmin\ upgrade\ REPOS\_PATH}

\subsection{Description}

Upgrade the repository located at \textless REPOS\_PATH\textgreater{} to
the latest supported schema version.

This functionality is provided as a convenience for repository
administrators who wish to make use of new Subversion functionality
without having to undertake a potentially costly full repository dump
and load operation. As such, the upgrade performs only the minimum
amount of work needed to accomplish this while still maintaining the
integrity of the repository. While a dump and subsequent load guarantee
the most optimized repository state, \texttt{svnadmin\ upgrade} does
not.

You should \emph{always} back up your repository before upgrading.

\subsection{Options}

None

\subsection{Examples}

Upgrade the repository at path \texttt{/var/repos/svn}:

\begin{verbatim}
$ svnadmin upgrade /var/repos/svn
Repository lock acquired.
Please wait; upgrading the repository may take some time...

Upgrade completed.
\end{verbatim}

\section{svnadmin verify}\label{svn.ref.svnadmin.c.verify}

\emph{Verify the data stored in the repository.}

\textbf{Synopsis}

\texttt{svnadmin\ verify\ REPOS\_PATH}

\subsection{Description}

Run this command if you wish to verify the integrity of your repository.
This basically iterates through all revisions in the repository by
internally dumping all revisions and discarding the
output---it\textquotesingle s a good idea to run this on a regular basis
to guard against latent hard disk failures and ``bitrot.'' If this
command fails---which it will do at the first sign of a problem---that
means your repository has at least one corrupted revision, and you
should restore the corrupted revision from a backup (you did make a
backup, didn\textquotesingle t you?).

\subsection{Options}

\begin{verbatim}
--check-normalization
--keep-going
--memory-cache-size (-M) ARG
--metadata-only
--quiet (-q)
--revision (-r) ARG
--transaction (-t) ARG
\end{verbatim}

\subsection{Examples}

Verify a hung repository:

\begin{verbatim}
$ svnadmin verify /var/svn/repos/ 
* Verified revision 1729.
\end{verbatim}

